% !TeX root = ../main.tex
\chapter{Methodology}
\label{chap:methodology}

This study examines CD4AI by comparing it with four cases: Divination, an
AI-enabled system extended in this work to implement the pattern, and three
tools whose documentation describes related workflows. This chapter describes
how the pattern was formulated (Section~\ref{sec:method-pattern}), how the
cases were selected (Section~\ref{sec:method-cases}), and what evidence was
collected about them (Section~\ref{sec:method-evidence}). It then describes
how each case was analyzed (Section~\ref{sec:method-analysis}), how the
findings were used to assess the pattern (Section~\ref{sec:method-refinement}),
and how AI tools were used in this work (Section~\ref{sec:method-ai}).

\section{Pattern Formulation}
\label{sec:method-pattern}

The idea for CD4AI arose from the author's professional work on AI-enabled
systems. Its initial version connected three stages in a feedback loop. A test
stage in the delivery pipeline allowed deployment only when the regression
suite passed. Monitoring collected data from the application in production.
Curation selected the executions that showed problems, which the team then
used to create fixes and new tests.

An early version of the pattern was presented at the MiniPLoP workshop held at
XP~2026, in Brazil~\parencite{Iasco:CD4AI:2026}. The discussion did not change
the pattern, but it produced recommendations on how to present it, as well as
the suggestion to examine MLflow (Section~\ref{sec:method-cases}).

The description was revised while Chapter~\ref{chap:cd4ai} was written and the
cases were studied; both activities contributed to the revisions. The
three-stage loop remained the same. The main change concerned how evaluation
results inform delivery. The initial version prescribed a single rule: the
whole regression suite had to pass. The revised description requires that
results be interpreted against evidence of previous behavior and leaves the
blocking rule to the team's delivery policy. It describes options for supplying
that evidence and applying the policy, such as identifying cases that track
known limitations, comparing results against one or more reference versions or
stored reference runs, and using repeated executions or statistical tests when
results vary between runs. In particular, the illustrative stored-history
statistical comparison in Section~\ref{sec:cd4ai-result-interpretation} was
inspired by the author's professional practice. Other revisions clarified the
stages; for example, a monitoring signal does not confirm a defect, and curation
prepares evaluation cases, while fixing the application belongs to development.

Chapter~\ref{chap:cd4ai} presents the current version of the pattern,
including its context, problem, forces, solution, and a worked example.

\section{Case Selection}
\label{sec:method-cases}

The study examines four cases of two kinds: an application implementation and
three external tool workflows. Divination is an AI-enabled system whose
evaluation, monitoring, and curation infrastructure was implemented in this
work, and its source code is publicly available. The other three---DeepEval
with Confident AI, Braintrust, and MLflow---are platforms and frameworks for
developing AI-enabled systems, examined through their documented workflows and
supporting capabilities. The cases are considered individually before their
similarities and differences are compared.

At the beginning of the project, available tools were reviewed for workflows
that appeared to share the pattern's structure. The three tools were selected
because their workflows extend beyond running evaluations to observing use,
reviewing interactions, and preparing evaluation cases from what is observed.
Where an open-source framework relies on a companion platform to cover this
cycle, both are examined together: DeepEval runs evaluations, while Confident
AI supports monitoring, human review, and experiment comparison. MLflow's own
guides describe the cycle, and a Databricks guide is used as a supplementary
source for monitoring live traffic. Pydantic Evals, from the Pydantic AI
project, was considered as an evaluation library but not selected, because it
focuses on running evaluations; Logfire, Pydantic's observability platform, was
not examined.

Braintrust and DeepEval were already known to the author as evaluation tools.
MLflow was suggested by an attendee at the MiniPLoP presentation as a possible
match for the pattern. Other platforms with a similar scope, such as
LangSmith, were not examined. Selection was therefore exploratory, guided by
prior knowledge of the tools and a recommendation, rather than a systematic
survey of the available tools. The cases support an examination of these
workflows, not an estimate of how widespread the pattern is across AI-enabled
systems.

\subsection{Divination}
\label{sec:method-divination}

Divination is a conversational assistant that answers questions about
Dungeons \& Dragons using retrieval-augmented generation. It was developed by
Luiz Carlos Costa da Silva as a capstone project~\parencite{Silva:Divination:2025}.
After the pattern was formulated, this work added infrastructure to implement
it: an evaluation pipeline with single-turn and conversational evaluations;
monitoring that stores the application's interactions in a database and flags
possible problems; a page for reviewing interactions; and a CI runner that
executes the evaluations, compares results against a stored reference baseline,
and applies a delivery policy.

The implementation avoids depending on a particular platform. The CI runner
can run on a developer's machine, in a Git hook, or on a CI server; its GitHub
Actions workflow is an optional adapter~\parencite{Divination:CI}. Implementing
the stages directly keeps them visible in the source code, where they can be
studied independently of any platform's features. The evaluations use
DeepEval, which was chosen in order to study it. Another evaluation library
could serve the same purpose, although the simulated conversations would
require an equivalent feature or custom code. Divination thus shows the
pattern implemented with general-purpose components, while the other cases
show how dedicated platforms support related workflows.

Because the author implemented this infrastructure to represent CD4AI, its
correspondence with the pattern illustrates one implementation. It does not
provide independent evidence that the same structure arose in another team's
practice.

\section{Evidence Collection}
\label{sec:method-evidence}

For the three tools, the evidence consists of their official documentation:
guides describing workflows and reference pages describing specific
capabilities. Because online pages change, the reviewed pages were archived
when they were accessed, and each reference gives its access date. For
Divination, the evidence consists of its source code and project
documentation at a fixed revision, so later changes to the project do not
alter the material analyzed.

% TODO: If the case repository is public before submission, describe it here
% (archived source snapshots, figure sources) and give its URL.

The tools were also tried in a test setting to understand their features,
and DeepEval was used in Divination's evaluations. None of them was used in a
production environment for this study. Divination's interface examples use
synthetic data and do not constitute an observed deployment.
% TODO: Describe any further hands-on testing done before submission.

Documentation can establish a recommended workflow or an available mechanism.
Source inspection can establish how a mechanism is implemented. Neither
establishes that a team completed the cycle or achieved a quality improvement.
Documentation may also be inaccurate or outdated, as discussed in
Section~\ref{sec:discussion-threats}.

\section{Analysis Procedure}
\label{sec:method-analysis}

The external projects do not refer to CD4AI. Mapping their documented workflows
to its stages is an interpretation made in this study, not a claim that they
adopted the pattern.

For each case, the documentation was first reviewed to determine whether it
describes a workflow resembling the pattern's feedback loop. The features it
describes were then mapped to the monitoring, curation, and evaluation stages,
and the connections between the stages were checked: whether observed
interactions can become evaluation cases, and whether those cases are evaluated
on later changes. Finally, the analysis examined how evaluation results inform
delivery decisions. AI assistants supported this review
(Section~\ref{sec:method-ai}).

Each case is examined for two aspects separately: whether it supports the
feedback loop connecting evaluation, monitoring, and curation (elements~1--3,
addressing RQ1), and how it interprets evaluation results and connects them to
delivery decisions, particularly when distinguishing continuing known
limitations from potential regressions (elements~4 and~5, addressing RQ2). A
case may support the feedback loop without specifying how evaluation results
govern delivery.

In assessing delivery mechanisms, the analysis distinguishes technical
capabilities from team policy. A documented assertion or failed CI job
establishes a mechanism for rejecting an evaluation result, whereas whether
that result prevents integration or deployment also depends on the surrounding
policy and its enforcement. Similarly, support for version comparison or case
metadata does not by itself establish how a team manages known limitations.
When the reviewed sources leave these matters unspecified, the case identifies
an evidence gap rather than inferring that the behavior is present or absent.

To apply this procedure consistently in Chapter~\ref{chap:case-studies}, each
case study follows a four-part structure:
\begin{description}
  \item[Context and Evidence.] Introduces the system or platform, defines the
    scope of the workflow examined, and identifies the primary sources
    inspected.
  \item[The Feedback Loop.] Traces how observations from use move through
    monitoring, curation, and evaluation, and identifies which transitions are
    automated or manual (addressing RQ1).
  \item[Result Interpretation and Delivery Decisions.] Examines how evaluation
    results are compared across versions, how variability and known limitations
    are handled, and how outcomes connect to CI checks or release decisions
    (addressing RQ2).
  \item[Limitations.] Identifies what the implementation or documented workflow
    leaves unresolved, what policy decisions remain with the team, and where
    the available evidence is incomplete.
\end{description}

\section{Assessment of the Pattern}
\label{sec:method-refinement}

The assessment brings together the case findings to address RQ3 by identifying
shared connections, implementation differences, and gaps in the available
evidence. A correspondence supports the use of CD4AI to describe the examined
workflow; it does not establish that adopting the pattern improves delivery
outcomes.

The findings may suggest changes to the pattern's description.
Chapter~\ref{chap:discussion} presents these implications and distinguishes
evidence about the examined cases from claims about other workflows.

\section{Use of AI Tools}
\label{sec:method-ai}

% TODO (Step 8): Confirm that every cited source and technical claim has been
% verified against the archived snapshots and source code before submission.
AI assistants (Claude Code, OpenAI Codex, and Google Antigravity) were used
during this project to support coding, documentation review, and writing. They
assisted in implementing Divination's evaluation, monitoring, and curation
infrastructure, exploring and summarizing tool documentation during the case
analysis, checking consistency across chapters, drafting portions of the
thesis, and revising the English prose for grammar, clarity, and style. The
author selected the cases and primary sources and is responsible for verifying
the cited documentation and source code, as well as for the implementation,
analysis, and final text.
